% 本文件由 LaTeX 排版 Agent 根据有效 translation.json 自动生成。
% author=Gregory Thaumaturgus; work_id=0603; title=教规书信（圣额我略·显行灵迹者）

\chapter{教规书信（圣额我略·显行灵迹者）}

% structure: p0001 source=h1 target=omitted reason=page-title-already-rendered-by-parent

% structure: p0002 source=h2 target=section reason=relative-heading-rank; base=h2

\section{第一规}

至圣之父，肉类并不成为我们的负担，倘若被俘者吃了征服者摆放在他们面前的食物，尤其是因为来自各方的报告一致说，进犯我们地区的蛮族并未向偶像献祭。因为宗徒说：\BibleQuote{食物是为肚腹，肚腹是为食物；但天主必要使这两样都毁坏。}\BibleRef{格前 6:13} 而且，洁净一切食物的救主也说：\BibleQuote{不是进入人内的使人污秽，而是从人出来的才使人污秽。}\BibleRef{玛 15:11} 这也适用于被蛮族玷污、凌辱了她们身体的被俘妇女。但是，如果这类人先前的生活证明她曾如经上所记载的那样，\BibleQuote{随从淫乱者的眼目}，那么淫乱的习惯显然在被俘期间也值得怀疑。这样的人不宜轻易与人在祈祷中相通。然而，如果有人曾生活在极其贞洁之中，并在过去显示出纯洁且毫无嫌疑的生活方式，而如今因不得已的强迫而陷于放纵，我们有一个榜样可作指引——即《申命纪》中的少女，有人在田间遇见她，强行与她交合。他说：\BibleQuote{对那少女，你不可做什么；在少女身上没有该死的罪，因为这事如同人起来攻击邻舍，将他杀了一样。少女喊叫，却无人救她。}\BibleRef{申 22:26}\par

% structure: p0004 source=h2 target=section reason=relative-heading-rank; base=h2

\section{第二规}

贪财是一项大恶；在一封信中不可能列出所有经文，其中不仅宣告抢劫是可怕且可憎的事，而且一般来说，贪心以及为了满足卑劣的获利欲望而干预他人之物的倾向也是如此。所有这种精神的人都从天主的教会中被逐出。但是，在入侵期间，在如此悲惨的忧伤和痛苦的哀号中，竟有人胆大妄为，将给所有人带来毁灭的危机视为自己牟利的时机，这只能说明他们是亵渎神明、为天主所憎恶、且无可比拟的邪恶之人。故此，似乎应当将这样的人逐出教会，以免天主的忿怒临到全体人民，尤其是那些身居要职却疏于查问的人。因为我害怕，正如经上所说，邪恶之人会连同自己一起毁灭正义之人。经上说：\BibleQuote{因为淫乱和贪财\Emph{是事}，因这些事，天主的忿怒才临到悖逆之子。所以你们不要与他们同伙。从前你们是黑暗，如今在主内却是光明，行事为人就当像光明的子女（光明的果子在于一切良善、公义和诚实），察验何为主所喜悦的事。不要与黑暗无益的工作有分，倒要责备它们；因为他们暗中所做的事，连提起来也是可耻的。但一切被责备的事，都被光显明出来。}\BibleRef{弗 5:5} 宗徒这样说了。但如果某些人，在和平时期已因先前的贪财而受了应得的惩罚，如今在患难时刻（即在蛮族入侵的祸患中）又转而沉溺于贪财，并从那些被完全掠夺、或被俘、或被杀害的人的血和毁灭中获利，那么还能期望什么呢？无非是那些如此热切争夺贪财的人，为自己和全体人民积蓄忿怒。\par

% structure: p0006 source=h2 target=section reason=relative-heading-rank; base=h2

\section{第三规}

看，匝塔的儿子阿干\EditorialNote{《若苏厄书》第七章}岂非在当灭之物上犯了罪，那时祸患不是临到了以色列全会众吗？他一人独犯其罪，但因他的罪而死的却不止他一人。我们也应将此时一切带获利性的东西，凡是我们并非正当地拥有、而是严格属于他人的财产，都算作当灭之物。因为阿干确实取了掠物；而现今这些人也取了掠物。但他取的是敌人的东西；而这些人现今取的是弟兄的东西，并以致命的获利来肥己。\par

% structure: p0008 source=h2 target=section reason=relative-heading-rank; base=h2

\section{第四规}

任何人都不要自欺，也不要以找到这样东西为借口。因为即使人找到任何东西，也不得藉此肥己。因为《申命纪》说：\BibleQuote{你若看见你兄弟的牛或羊在路上失迷，不可佯为不见，总要把它牵回来交给你的兄弟。你兄弟若离你远，或是你不认识他，就要牵到你家去，留在你那里，等你兄弟来寻找，就还给他。他的驴你也要照样行，他的衣服你也要照样行，凡你兄弟遗失的一切东西，你若找到，都要照样行。}\BibleRef{申 22:1} 这是《申命纪》的话。而在《出谷纪》中，不仅涉及找到朋友之物，也涉及找到仇敌之物，经上说：\BibleQuote{总要把它牵回你主人的家。}\BibleRef{出 23:4} 如果即使在和平时期，当他人安逸、舒适、无忧无虑地生活时，也不得藉他人（无论是弟兄还是仇敌）之物肥己，那么当人遭遇逆境、逃离敌人、因形势所迫而不得不放弃自己的财产时，岂非更应如此？\par

% structure: p0010 source=h2 target=section reason=relative-heading-rank; base=h2

\section{第五规}

但另一些人自欺欺人，以为可以保留他人之物，作为自己失物的补偿。诚然，正如\Person{Boradi}人和\Person{Goth}人给他们带来战祸，他们自己也成了他人的\Person{Boradi}人和\Person{Goth}人。因此，我们派了我们的弟兄、我们的老同伴\Person{Euphrosynus}到你们那里，目的就是让他按照我们这里的模式与你们相处，教导你们应接纳哪些人的控告，以及应从你们的祈祷中排除哪些人。\par

% structure: p0012 source=h2 target=section reason=relative-heading-rank; base=h2

\section{第六条规则}

此外，有人向我们报告，在你们的地区发生了一件确实令人难以置信的事；如果确有其事，那完全是出于不信者和不虔敬之人，是那些连主的名都不知道的人所为：有些人竟残忍和不人道到如此地步，以致强行扣留某些逃跑的俘虏。请立即派遣专员到乡间去，免得天主的雷霆确实落在行此事的人身上。\par

% structure: p0014 source=h2 target=section reason=relative-heading-rank; base=h2

\section{第七条规则}

至于那些曾被编入蛮族之中、以俘虏身份随同他们入侵的人，以及那些忘记自己来自\Place{本都}、忘记自己是基督徒，竟变得如此彻底蛮化，甚至用绞刑或勒杀杀害自己的同族，并为原本不识路的蛮族指明道路或房屋的人，你们必须禁止这样的人甚至在公开集会中作为听道者，直到诸圣在会议中，并借着在他们之前的圣神，对此作出共同决定为止。\par

% structure: p0016 source=h2 target=section reason=relative-heading-rank; base=h2

\section{第八条规则}

那些胆敢侵入他人房屋的人，如果已经受审并被定罪，就不应被认为配得上在公开集会中作听道者。但如果他们自行坦白并归还财物，则应被置于悔罪者的行列。\par

% structure: p0018 source=h2 target=section reason=relative-heading-rank; base=h2

\section{第九条规则}

那些在田野里或自己家中找到蛮族遗落之物的人，如果已经受审并被定罪，应归入同一类悔罪者。但如果他们自行坦白并归还财物，则应被认为配得上祈祷的特权。\par

% structure: p0020 source=h2 target=section reason=relative-heading-rank; base=h2

\section{第十条规则}

遵守诫命的人应当毫无卑劣贪心地遵守它，既不要求报酬、赏金、费用，也不要求任何其他被称为谢礼的东西。\par

% structure: p0022 source=h2 target=section reason=relative-heading-rank; base=h2

\section{第十一条规则}

哀哭发生在圣堂大门之外；罪人站在那里，应当恳求进入的信友为他祈祷。听道则发生在大门内的门廊里，罪人应当站在那里，直到慕道者\Emph{离开}，然后他才出去。因为经上说：让他聆听圣经和道理，然后被赶出去，被认为不配享有祈祷的特权。臣服是指一个人站在圣殿的大门内，同慕道者一起出去。恢复是指一个人与信友共处，而不与慕道者一同出去；最后是参与神圣礼仪。\par

\SourceRecord{Translated by S.D.F. Salmond.；Ante-Nicene Fathers；Vol. 6.；Edited by Alexander Roberts, James Donaldson, and A. Cleveland Coxe.；Buffalo, NY: Christian Literature Publishing Co.,；1886.；<http://www.newadvent.org/fathers/0603.htm>.}
